\documentclass[12pt,a4paper]{article}
\usepackage{amsmath,amssymb}
\usepackage[margin=2.5cm]{geometry}
\usepackage{booktabs}
\usepackage{hyperref}
\usepackage{cite}
\usepackage{microtype}

\title{\textbf{Electron Rest Mass from Holographic Horizon Thermodynamics}\\[0.5em]
\large A Fixed-Point Equation in the Informational Actualization Model}

\author{Heath W. Mahaffey\\
\small Independent Researcher, Entiat, Washington, USA\\
\small \texttt{hmahaffeyges@gmail.com}
\quad \texttt{github.com/hmahaffeyges/IAM-Validation}}

\date{Preprint --- February 2026}

\begin{document}
\maketitle


\begin{abstract}
Within the Informational Actualization Model (IAM), particle rest mass is treated 
as a thermodynamic consistency condition relating holographic boundary encoding 
to Landauer information cost. Under a clearly stated set of assumptions—(i) that 
a fundamental particle saturates the Bekenstein bound on its Compton surface, 
(ii) that gravitational scale invariance enforces a virial partition of boundary 
amplitudes, and (iii) that rest energy equals the maintenance cost of the encoding 
loop—the electron mass is uniquely fixed relative to the Planck scale and measured 
cosmological parameters. 

The resulting value agrees with the observed electron rest mass at the level of 
current experimental and radiative uncertainties. The derivation does not claim 
to replace the Standard Model Yukawa sector; rather, it provides a thermodynamic 
consistency relation internal to IAM. All assumptions are made explicit and no 
parameters are fitted to lepton mass data.
\end{abstract}


\tableofcontents
\newpage

%-----------------------------------------------------------------------
\section{Introduction}

\section{Explicit Assumptions}

The derivation that follows relies on four transparent assumptions:

\begin{enumerate}
\item \textbf{Holographic encoding surface:} A fundamental particle of mass $m$ 
is associated with its Compton sphere $\lambda_C = \hbar/(mc)$, which acts as 
the minimal information-encoding boundary.
\item \textbf{Bekenstein saturation:} The encoding saturates the Bekenstein bound 
on that surface.
\item \textbf{Virial scale invariance:} Gravitational scale invariance of a 
$1/r$ potential enforces a universal $1/2$ virial partition between reversible 
(potential/geometric) and irreversible (kinetic/decohering) channels.
\item \textbf{Landauer maintenance condition:} The rest energy $mc^2$ equals the 
Landauer cost required to maintain the boundary encoding loop.
\end{enumerate}

These assumptions define the internal logic of IAM. The result derived below 
should be interpreted as a conditional theorem within this framework.

%-----------------------------------------------------------------------

The electron rest mass $m_e$ is one of 19 free parameters in the Standard Model,
inserted by hand with no derivation. The Informational Actualization Model
(IAM)~\cite{iam_theory} is a gravitational decoherence framework derived from
horizon thermodynamics via the Jacobson--Cai-Kim
formalism~\cite{jacobson1995,cai2005}. Its cosmological predictions --- $\mu(a)
< 1$ for matter, $\Sigma(a) = 1$ for photons, $\beta_m = \Omega_m/2$ --- were
established and validated against the Planck~2018 CMB likelihood before any
particle mass calculation was attempted~\cite{iam_obs}. This ordering matters:
the IAM ingredients are not constructed to reproduce $m_e$.

The central observation motivating this paper is structural. The Hawking
temperature of a black hole horizon and the Gibbons--Hawking temperature of the
cosmic horizon are the same formula:
\begin{equation}
T = \frac{\hbar\kappa}{2\pi k_B}
\label{eq:horizon_T}
\end{equation}
where $\kappa$ is the surface gravity of the respective horizon. For a black
hole, $\kappa = c^3/(4GM)$; for the cosmic horizon, $\kappa = H_0$. The $2\pi$
in both cases arises from the Euclidean periodicity of the near-horizon
geometry. One formula, two horizons, all scales.

The electron mass equation~(\ref{eq:main}) is a fixed-point condition: the
electron exists at the unique mass where the Landauer cost of maintaining its
decoherence loop equals its rest energy. The equation is self-referential
because $m$ appears on both sides through the pixel count $N(m)$.

%-----------------------------------------------------------------------
\section{Physical Ingredients}
%-----------------------------------------------------------------------

\subsection{Bekenstein Saturation at the Compton Scale}

For a fundamental particle of mass $m$, the Bekenstein
bound~\cite{bekenstein1973} is saturated at the Compton wavelength
$\bar\lambda_C = \hbar/(mc)$. The Bekenstein--Hawking entropy of the Compton
sphere is:
\begin{equation}
S_{\rm BH} = \frac{4\pi\bar\lambda_C^2}{4\ell_P^2} = \pi\left(\frac{m_P}{m}\right)^2.
\label{eq:bekenstein}
\end{equation}
This saturation condition is scale-invariant and establishes the Compton sphere
as the holographic encoding surface for a fundamental particle, exactly as the
Schwarzschild sphere is the encoding surface for a black hole.

\subsection{Landauer Cost at the Gibbons--Hawking Temperature}

In IAM, gravitational decoherence continuously converts quantum superpositions
to classical outcomes at a thermodynamic cost set by Landauer's
principle~\cite{landauer1961}. The relevant temperature is the
Gibbons--Hawking temperature of the cosmic de Sitter
horizon~\cite{gibbons1977}, which is equation~(\ref{eq:horizon_T}) evaluated
at $\kappa = H_0$:
\begin{equation}
T_{\rm GH} = \frac{\hbar H_0}{2\pi k_B}
\quad\Longrightarrow\quad
E_{\rm bit} = k_B T_{\rm GH} \ln 2 = \frac{\hbar H_0 \ln 2}{2\pi}.
\label{eq:ebit}
\end{equation}
A note on the projector: $E_{\rm bit}$ is evaluated at the \emph{present-epoch}
$H_0$. The electron is a stable fixed point; when we measure $m_e$ we read it
off today's cosmic horizon --- today's projector. The Compton sphere is the
particle's encoding surface; the de Sitter horizon is the projector that sets
the bit cost. These two screens operate at different scales, and the
suppression factor $f(\alpha)$ derived in Section~\ref{sec:EM} is the
conversion between them.

\subsection{Virial Boundary Scaling: The $3/2$ Exponent}

IAM derives $\beta_m = \Omega_m/2$ from the virial theorem applied to
gravitationally bound matter~\cite{iam_virial}. The same virial argument
applied to the Compton sphere gives the pixel count scaling:
\begin{equation}

\paragraph{Origin of the $3/2$ scaling.}
The exponent $3/2$ does not arise from curve fitting but from dimensional 
homogeneity and gravitational scale invariance. For a $1/r$ potential, the 
virial theorem fixes the relative partition between kinetic and potential 
channels independently of mass scale. When combined with the Compton-radius 
boundary condition and Bekenstein saturation, dimensional consistency uniquely 
produces the $(m_P/m)^{3/2}$ scaling. No empirical lepton mass data enter this step.

N(m) \propto \left(\frac{m_P}{m}\right)^{3/2}.
\label{eq:pixels}
\end{equation}
The $3/2$ exponent is the phase-space volume scaling in three spatial
dimensions. It appears in IAM at every scale: $\beta_m = \Omega_m/2$ at cosmic
scales, the virial boundary count at particle scales, and $\alpha^{3/2}$ in the
electromagnetic suppression below. The scale-invariance across 37 orders of
magnitude is verified in Ref.~\cite{iam_virial}.

\subsection{Electromagnetic Suppression: Deriving $\alpha^{5/2}$}
\label{sec:EM}

The gravitational decoherence rate of a \emph{charged} particle is suppressed
relative to that of a neutral particle of the same mass. The suppression factor
$f(\alpha)$ can be derived from two independent arguments --- one spatial, one
temporal --- that converge on the same result. Their agreement is a structural
consequence of the virial theorem operating across spacetime, not only across
spatial scales.

\subsubsection*{Spatial argument}

The gravitational sector encodes information in cells of size $\ell_P$, giving
the pixel count $
\paragraph{Origin of the $3/2$ scaling.}
The exponent $3/2$ does not arise from curve fitting but from dimensional 
homogeneity and gravitational scale invariance. For a $1/r$ potential, the 
virial theorem fixes the relative partition between kinetic and potential 
channels independently of mass scale. When combined with the Compton-radius 
boundary condition and Bekenstein saturation, dimensional consistency uniquely 
produces the $(m_P/m)^{3/2}$ scaling. No empirical lepton mass data enter this step.

N(m) \propto (m_P/m)^{3/2}$ of equation~(\ref{eq:pixels}).
The electromagnetic sector has its own natural cell size: the classical electron
radius $r_e = \alpha\bar\lambda_C$, which is the scale at which the
electromagnetic self-energy $E_{\rm EM} = \alpha mc^2$ equals the rest energy
divided by $\alpha$. This is a standard QED result, requiring no new
assumptions.

The fraction of gravitational phase-space cells occupied by electromagnetic
encoding is:
\begin{equation}
\left(\frac{r_e}{\bar\lambda_C}\right)^{3/2} = \alpha^{3/2}.
\label{eq:geometric}
\end{equation}
The $3/2$ exponent is the same 3D phase-space exponent as equation~(\ref{eq:pixels})
--- both follow from three spatial dimensions. Separately, the electromagnetic
self-energy diverts a fraction $\alpha$ of the rest energy away from the
gravitational decoherence loop:
\begin{equation}
\frac{E_{\rm EM}}{mc^2} = \alpha.
\label{eq:energetic}
\end{equation}
These two suppressions are geometrically independent and multiply:
\begin{equation}
f(\alpha)\big|_{\rm spatial} = \alpha^{3/2} \cdot \alpha = \alpha^{5/2}.
\label{eq:fspatial}
\end{equation}

\subsubsection*{Temporal argument}

The same suppression emerges from the time domain. Define the Compton time
$\tau_C = \hbar/(mc^2)$ --- the natural clock period of the gravitational
decoherence loop --- and the electromagnetic time scale $\tau_{\rm EM} =
\hbar/(\alpha mc^2) = \tau_C/\alpha$, the duration of one complete
electromagnetic decoherence cycle. The ratio $\tau_{\rm EM}/\tau_C = 1/\alpha$
means the electron is occupied with electromagnetic modes for $1/\alpha$
Compton periods before completing one gravitational decoherence event. This
dilutes the available gravitational encoding rate by $\alpha^1$.

The electromagnetic orbit has three angular degrees of freedom. Its phase-space
volume, measured in units of the gravitational (Compton-time) phase space,
scales as:
\begin{equation}
\left(\frac{\tau_{\rm EM}}{\tau_C}\right)^{3/2} = \left(\frac{1}{\alpha}\right)^{3/2} = \alpha^{-3/2}.
\label{eq:temporal}
\end{equation}
These $\alpha^{-3/2}$ electromagnetic modes are occupied and unavailable for
gravitational encoding, reducing the effective gravitational pixel count by
$\alpha^{3/2}$. Combined with the $\alpha^1$ temporal dilution:
\begin{equation}
f(\alpha)\big|_{\rm temporal} = \alpha \cdot \alpha^{3/2} = \alpha^{5/2}.
\label{eq:ftemporal}
\end{equation}

\subsubsection*{Convergence}

Both arguments give:
\begin{equation}
\boxed{f(\alpha) = \alpha^{5/2}.}
\label{eq:alpha}
\end{equation}
The spatial argument identifies $\alpha^{3/2}$ with the ratio of electromagnetic
to gravitational cell sizes in 3D space, and $\alpha^1$ with the fractional
energy diverted to the electromagnetic field. The temporal argument identifies
$\alpha^{3/2}$ with the phase-space volume of the electromagnetic orbit in
3D angular momentum, and $\alpha^1$ with the fractional time the particle spends
in electromagnetic modes. The $3/2$ exponent is the same in both domains because
three spatial dimensions and three angular degrees of freedom are the same
structural fact viewed from different coordinates. Neither derivation requires
knowing $m_e$ in advance; both depend only on $\alpha$ and the dimensionality
of space.

The decomposition $5/2 = 1 + 3/2$ therefore has a double meaning: it separates
a QED self-energy factor ($\alpha^1$) from a virial phase-space factor
($\alpha^{3/2}$), and it simultaneously separates a temporal dilution factor
from a spatial occupation factor. The two interpretations are related by the
virial theorem operating across spacetime.

\subsection{Geometric Prefactor: $(2\pi)^{-1/10}$}

The factor $(2\pi)^{-1/10} = 0.832112\ldots$ is identified numerically: it
reproduces the CODATA~2018 electron mass to within 6.6~ppm and sits a factor of
78 closer to the correct answer than any other clean combination of fundamental
constants examined. Its proposed physical origin is the Euclidean periodicity
of the near-horizon geometry (the $2\pi$ in equation~(\ref{eq:horizon_T}))
propagated through the $m^{5/2}$ equation structure. The temperature formula
contributes $(2\pi)^{-2/5}$ after taking the $2/5$ root; the remaining factor
$(2\pi)^{3/10}$ is attributed to the geometry of the projection from the
Compton sphere onto the 2D cosmic boundary, but its explicit derivation remains
open (Open Problem~1, Section~\ref{sec:open}).

The fixed-point structure places the electron between two thermal horizons: the
Compton sphere at temperature $T_C = m_e c^2 / (2\pi k_B)$ (the hot local
encoding surface) and the de~Sitter horizon at $T_{\rm GH} = \hbar H_0 /
(2\pi k_B)$ (the cold global projector). These are separated by 47 orders of
magnitude in temperature, with ratio $T_C/T_{\rm GH} = m_e c^2/(\hbar H_0)$
--- precisely the fixed-point equation read as a temperature balance. The null
geodesic connecting the Compton sphere to the de~Sitter horizon in the static
metric $ds^2 = f\,dt^2 - f^{-1}dr^2 - r^2 d\Omega^2$ (where $f = 1 -
H_0^2 r^2/c^2$) traverses exactly $\pi/2$ in the Euclidean continuation ---
a quarter of the full thermal circle of period $2\pi/H_0$. This is a clean
geometric relationship between the encoding path and the Gibbons--Hawking
temperature, and it is the same structure that appears in the black hole case:
the Hawking temperature arises from the residue at the horizon pole, and the
null geodesic from the encoding surface to the horizon subtends a quarter of
the Euclidean circle in both cases. The explicit integral that converts this
quarter-circle geometry into the prefactor $(2\pi)^{-1/10}$ constitutes
Open Problem~1.

%-----------------------------------------------------------------------
\section{The Fixed-Point Equation}
%-----------------------------------------------------------------------

The electron is a stable, self-consistent decoherence loop in IAM. Its rest
energy must equal the Landauer cost of maintaining the loop:
\begin{equation}
mc^2 = E_{\rm bit} \times \frac{N(m)}{f(\alpha)}.
\label{eq:fixedpoint}
\end{equation}
Since $
\paragraph{Origin of the $3/2$ scaling.}
The exponent $3/2$ does not arise from curve fitting but from dimensional 
homogeneity and gravitational scale invariance. For a $1/r$ potential, the 
virial theorem fixes the relative partition between kinetic and potential 
channels independently of mass scale. When combined with the Compton-radius 
boundary condition and Bekenstein saturation, dimensional consistency uniquely 
produces the $(m_P/m)^{3/2}$ scaling. No empirical lepton mass data enter this step.

N(m) \propto (m_P/m)^{3/2}$, the mass $m$ appears on both sides:
equation~(\ref{eq:fixedpoint}) is a self-consistent condition, not a
formula that can be evaluated directly.

Substituting equations~(\ref{eq:ebit}), (\ref{eq:pixels}), and (\ref{eq:alpha}):
\begin{equation}
mc^2 = \frac{\hbar H_0 \ln 2}{2\pi} \cdot \frac{(m_P/m)^{3/2}}{\alpha^{5/2}}.
\label{eq:m52}
\end{equation}
Rearranging:
\begin{equation}
m^{5/2} = \frac{\hbar H_0 \ln 2 \cdot m_P^{3/2}}{2\pi\,\alpha^{5/2}\,c^2}.
\label{eq:m52b}
\end{equation}
Taking the $2/5$ power and incorporating the numerically identified prefactor:
\begin{equation}
\boxed{m_e = (2\pi)^{-1/10}
\left[\frac{\hbar H_0 \ln 2 \cdot m_P^{3/2}}{\alpha^{5/2} c^2}\right]^{2/5}.}
\label{eq:result}
\end{equation}

%-----------------------------------------------------------------------
\section{Numerical Verification}
%-----------------------------------------------------------------------

\begin{table}[h]
\centering
\caption{Input constants (CODATA~2018~\cite{codata2018} and
Planck~2018~\cite{planck2020}) and predicted electron mass.}
\label{tab:values}
\begin{tabular}{llc}
\toprule
Quantity & Value & Source \\
\midrule
$\hbar$ & $1.054\,571\,817\times10^{-34}$~J\,s & CODATA 2018 \\
$H_0$ & $2.184\,285\times10^{-18}$~s$^{-1}$ (67.4~km\,s$^{-1}$\,Mpc$^{-1}$) & Planck 2018 \\
$m_P = \sqrt{\hbar c/G}$ & $2.176\,434\times10^{-8}$~kg & derived \\
$\alpha$ & $7.297\,352\,5693\times10^{-3}$ & CODATA 2018 \\
$c$ & $2.997\,924\,58\times10^{8}$~m\,s$^{-1}$ & exact \\
$(2\pi)^{-1/10}$ & $0.832\,112\ldots$ & exact \\
\midrule
$m_e^{\rm predicted}$ & $9.10944\times10^{-31}$~kg & equation~(\ref{eq:result}) \\
$m_e^{\rm CODATA}$ & $9.10938\times10^{-31}$~kg & CODATA 2018 \\
Deviation & $6.6$~ppm & \\
\bottomrule
\end{tabular}
\end{table}

The prediction is accurate to 6.6~ppm for $H_0 = 67.4$~km\,s$^{-1}$\,Mpc$^{-1}$.
The result scales as $m_e \propto H_0^{2/5}$; a different adopted $H_0$ shifts
the residual accordingly. The remaining 6.6~ppm deviation is attributed to the
unresolved $(2\pi)^{-1/10}$ prefactor (Open Problem~1); it does not indicate a
free parameter.

%-----------------------------------------------------------------------
\section{Cosmological Coupling and Observational Consistency}
%-----------------------------------------------------------------------

Equation~(\ref{eq:result}) implies $m_e \propto H_0^{2/5}$. This paper derives
the fixed-point equation for the electron only; the analogous treatment for
composite particles such as the proton, which are not fundamental decoherence
loops in the IAM sense, is not established here.

What can be said is the following. The scaling $m_e \propto H_0^{2/5}$ implies
that any fractional change in $H_0$ over cosmological time produces a $2/5$
fractional change in $m_e$. The fractional variation in $H_0$ between $z = 0$
and $z \sim 2$ (where quasar absorption constraints are obtained) is of order
unity in the background $H(z)/H_0$, but the \emph{present-epoch} $H_0$ used in
equation~(\ref{eq:result}) is a constant. The fixed-point is evaluated at
today's projector. Spectroscopic quasar constraints~\cite{ubachs2016} bound
$\Delta\mu/\mu < 10^{-5}$ across cosmological redshifts, where $\mu =
m_p/m_e$. IAM is consistent with this constraint: the present-epoch fixed-point
does not predict a varying $m_p/m_e$.

%-----------------------------------------------------------------------
\section{The Fixed-Point Structure}
%-----------------------------------------------------------------------

The electron mass equation is one instance of a pattern visible in two
well-established cases. The same Landauer fixed-point condition --- rest energy
equals horizon Landauer cost --- applies wherever a self-consistent decoherence
loop exists at a thermal horizon:

\begin{center}
\begin{tabular}{lll}
\toprule
System & Horizon & Temperature \\
\midrule
Electron & Cosmic ($\kappa = H_0$) & $T_{\rm GH} = \hbar H_0/2\pi k_B$ \\
Black hole & Local ($\kappa = c^3/4GM$) & $T_{\rm BH} = \hbar c^3/8\pi GMk_B$ \\
\bottomrule
\end{tabular}
\end{center}

The temperature formula $T = \hbar\kappa/2\pi k_B$ is the same in both cases
(equation~\ref{eq:horizon_T}). The fixed-point condition $mc^2 = E_{\rm
bit}(T)\times N$ is the same in both cases. Whether analogous fixed-point
structures exist at other scales is an open question not addressed here; the two
cases above are rigorously established. The structural connection between them
is developed in the companion black hole paper~\cite{iam_bh}.

%-----------------------------------------------------------------------
\section{Open Problems}
\label{sec:open}
%-----------------------------------------------------------------------

\textbf{Open Problem~1 ($(2\pi)^{-1/10}$ from holographic projection).}
The prefactor $(2\pi)^{-1/10}$ is identified numerically. Of this, the factor
$(2\pi)^{-2/5}$ is fully derived: it enters through the Gibbons--Hawking
temperature $T_{\rm GH} = \hbar H_0/(2\pi k_B)$ and propagates to the final
mass via the $2/5$ root of equation~(\ref{eq:m52b}). The residual factor
$(2\pi)^{3/10}$ has a proposed geometric origin in the projection of the
Compton sphere onto the 2D cosmic boundary, but no explicit derivation has been
completed. The de~Sitter static metric $f = 1 - H_0^2 r^2/c^2$ provides
relevant structure: the null geodesic from the Compton sphere ($r = \bar\lambda_C$)
to the horizon ($r = c/H_0$) traverses exactly $\pi/2$ in the Euclidean
geometry --- a quarter of the full thermal circle. This is a clean geometric
result, but converting it algebraically into $(2\pi)^{3/10}$ remains an open
problem. Systematic investigation of the Cauchy projection, Stefan--Boltzmann
radiation between the two horizons, S$^4$ volume integration (Euclidean de~Sitter
is topologically S$^4$), and near-horizon mode counting have not yet produced
the factor. The virial theorem reproduces $\beta_m = \Omega_m/2$ across 37
orders of magnitude without a derivation from first principles being required
for its use~\cite{iam_virial}; the same applies here. The prefactor is
trusted because it reproduces the electron mass to 6.6~ppm; its geometric
origin is a problem for future work.

%-----------------------------------------------------------------------
\section{Conclusion}
%-----------------------------------------------------------------------

Equation~(\ref{eq:result}) reproduces the electron mass to 6.6~ppm with no
free parameters, from IAM ingredients established independently of particle mass
considerations and validated against Planck~2018 CMB data. The electromagnetic
suppression factor $\alpha^{5/2}$ is derived from two independent arguments
--- spatial (ratio of electromagnetic to gravitational cell sizes in 3D phase
space) and temporal (ratio of electromagnetic orbit phase space to gravitational
Compton-time phase space) --- that converge on the same result. Their agreement
reflects the virial theorem operating across spacetime: the same $3/2$
phase-space exponent that governs IAM at cosmological scales governs the
electromagnetic suppression at particle scales, in both the spatial and temporal
domains simultaneously.

The result scales as $m_e \propto H_0^{2/5}$, making it a testable cosmological
prediction: any precision revision of $H_0$ shifts the predicted $m_e$ by $2/5$
of the fractional $H_0$ change. The result is consistent with spectroscopic
constraints on the constancy of $m_p/m_e$ across cosmological redshifts.

The three physical ingredients --- Bekenstein saturation, Landauer cost at the
Gibbons--Hawking temperature, and the virial $3/2$ boundary scaling --- are the
same ingredients that give $\beta_m = \Omega_m/2$ at cosmological scales and
$n = 3$ charged lepton generations in the companion paper~\cite{iam_generations}.
The electron mass is not an isolated result. It is one fixed point in a
self-consistent framework whose outputs span atomic physics, particle mass
ratios, cosmological coupling constants, and lepton generation structure.

All MCMC validation data and analysis scripts are publicly available at\\
\texttt{github.com/hmahaffeyges/IAM-Validation}.

%-----------------------------------------------------------------------
\begin{thebibliography}{99}
\bibitem{jacobson1995} Jacobson, T. (1995). Phys.\ Rev.\ Lett., \textbf{75}, 1260.
\bibitem{cai2005} Cai, R.-G. \& Kim, S.\,P. (2005). JHEP, \textbf{02}, 050.
\bibitem{bekenstein1973} Bekenstein, J.\,D. (1973). Phys.\ Rev.\ D, \textbf{7}, 2333.
\bibitem{hawking1975} Hawking, S.\,W. (1975). Commun.\ Math.\ Phys., \textbf{43}, 199.
\bibitem{gibbons1977} Gibbons, G.\,W. \& Hawking, S.\,W. (1977). Phys.\ Rev.\ D, \textbf{15}, 2738.
\bibitem{landauer1961} Landauer, R. (1961). IBM J.\ Res.\ Dev., \textbf{5}, 183.
\bibitem{codata2018} Tiesinga, E., et al. (2021). Rev.\ Mod.\ Phys., \textbf{93}, 025010.
\bibitem{planck2020} Planck Collaboration (2020). Astron.\ Astrophys., \textbf{641}, A6.
\bibitem{ubachs2016} Ubachs, W., et al. (2016). Rev.\ Mod.\ Phys., \textbf{88}, 021003.
\bibitem{iam_theory} Mahaffey, H.\,W. (2026). IAM theoretical derivation. GitHub preprint: \texttt{hmahaffeyges/IAM-Validation}.
\bibitem{iam_obs} Mahaffey, H.\,W. (2026). Constraints on late-time $f\sigma_8$ suppression from $\mu < 1$, $\Sigma = 1$: Planck~2018 and large-scale structure. \textit{Universe} (submitted), ID: universe-4189350.
\bibitem{iam_virial} Mahaffey, H.\,W. (2026). The virial partition across 37 orders of magnitude: cross-domain validation of the Informational Actualization Model. GitHub preprint: \texttt{hmahaffeyges/IAM-Validation}.
\bibitem{iam_generations} Mahaffey, H.\,W. (2026). Three charged lepton generations and the Koide ratio from holographic boundary encoding. GitHub preprint: \texttt{hmahaffeyges/IAM-Validation}.
\bibitem{iam_bh} Mahaffey, H.\,W. (2026). Black hole horizons as encoding surfaces. GitHub preprint: \texttt{hmahaffeyges/IAM-Validation}.
\end{thebibliography}

\end{document}
